The AI research community has adopted the word ``curiosity'' but means something quite specific and different by it.

\subsection{Artificial Curiosity as Intrinsic Motivation}\label{subsec:intrinsic-motivation}

In reinforcement learning, ``curiosity'' is operationalized as an \textbf{intrinsic reward signal} that motivates exploration in the absence of external rewards. The dominant frameworks include:

\begin{itemize}[leftmargin=*]
  \item \textbf{Prediction-error curiosity} \citep{schmidhuber1991,pathak2017}: The agent receives a reward proportional to how poorly it predicted the outcome of its action. Novelty equals prediction error equals reward signal. The agent seeks states that surprise it.

  \item \textbf{Count-based novelty} \citep{bellemare2016}: The agent is rewarded for visiting states it has not visited often. Rare equals interesting.

  \item \textbf{Random Network Distillation} \citep{burda2019}: A fixed random network provides a target; a learned network tries to match its outputs. The prediction error between them serves as a novelty signal.
\end{itemize}

These are genuinely useful for solving hard exploration problems---agents with artificial curiosity can learn to play Montezuma's Revenge, navigate mazes, and solve sparse-reward environments that defeat standard RL.

\textbf{The key observation:} In all these frameworks, ``curiosity'' is a \emph{functional mechanism}---a reward-shaping technique. The agent does not \emph{feel} an information gap. It receives a numerical bonus for visiting novel states. The word ``curiosity'' is used by analogy with the behavioral outcome (exploration of novelty), not the subjective experience.

\subsection{Curiosity in LLMs}\label{subsec:llm-curiosity}

Two recent papers directly evaluate whether LLMs exhibit something like curiosity.

\citet{wang2025} applied the 5DCR psychological curiosity scale to LLMs. Key findings:
\begin{itemize}[leftmargin=*]
  \item LLMs show \textbf{stronger information-seeking behavior} than typical human baselines.
  \item But they \textbf{make conservative choices in uncertain environments}---they prefer known-safe options over risky-novel ones.
  \item Curiosity-like behaviors \textbf{enhance reasoning and active learning} in LLMs.
  \item The authors conclude LLMs ``have the potential to exhibit curiosity similar to that of humans.''
\end{itemize}

\citet{borah2025} examined curiosity expression cross-culturally using their CUEST framework (measuring linguistic style, topic preferences, and social science grounding across 18 countries). Their findings are sobering:
\begin{itemize}[leftmargin=*]
  \item LLMs \textbf{compress cultural variation in curiosity by 63\%} (human standard deviation across countries: 0.0785; LLM standard deviation: 0.029). They produce a homogenized, Western-normative version of curiosity regardless of cultural persona.
  \item LLMs ask \textbf{inward-facing, socially neutral questions} while humans ask outward-facing, contrastive, cross-cultural ones. LLMs rarely contrast countries or ask causal/historical questions.
  \item \textbf{Smaller models with less RLHF outperform larger ones.} LLaMA-3-8b consistently beats GPT-5, Claude, and GPT-4o at matching human curiosity diversity. Safety training actively suppresses cultural variation.
  \item LLMs \textbf{reinforce cultural stereotypes}---they align best where human curiosity matches stereotypical expectations and break down when humans deviate from stereotypes.
  \item Fine-tuning (adapter-based with conversational scaffolding) can narrow the gap by up to 50\%, and curiosity-enabled models show 8+ percentage point gains on cultural reasoning benchmarks. But the underlying ``curiosity'' remains substantially a training data artifact.
\end{itemize}

\subsection{Linguistic Curiosity as a Distinct Channel}\label{subsec:linguistic-curiosity}

One finding from the developmental robotics work deserves special attention for LLMs. \citet{tinker2025} separated curiosity into modality-specific channels (vision, touch, proprioception, command voice, feedback voice). The critical discovery: including \emph{feedback voice curiosity}---where the robot is rewarded for encountering novel linguistic descriptions of its own actions---boosted generalization from approximately 75\% to approximately 85\%. The agent that wanted to \emph{hear new things said about what it did} learned compositional language faster than the agent that only sought novel sensory experiences.

This is provocative for LLMs. Language models operate almost entirely in the linguistic channel. If curiosity can be meaningfully instantiated through language alone---if wanting to encounter novel linguistic patterns is itself a form of curiosity---then the question of whether LLMs can be ``genuinely curious'' becomes harder to dismiss. The robot's linguistic curiosity channel was the smallest part of its architecture but produced disproportionate learning gains.

\subsection{The Emerging Hybrid Approach}\label{subsec:hybrid}

Recent work by \citet{quadros2025} combines LLM-derived intrinsic rewards with traditional novelty signals (VAE-based state novelty). The LLM provides semantic guidance (``this seems worth exploring because\ldots'') while the novelty detector provides the raw signal. This is interesting precisely because it acknowledges that neither mechanism alone does what human curiosity does.
